\chapter{The Zeros of \texorpdfstring{$\zeta$}{zeta}}\label{ch:zeros}

\begin{center}\itshape
One now finds indeed approximately this number of real roots within these limits, and it is very probable that all roots are real.\\
\upshape --- Riemann, \emph{On the number of primes less than a given magnitude} (1859)
\end{center}

Riemann's memoir is eight pages long. In those pages he extends $\zeta$ to the plane, proves the functional equation, and then makes three claims about the zeros in the strip $0<\operatorname{Re}s<1$.

His $\xi$ is a function of a real variable $t$. Set $s=\frac12+it$. The function he calls $\xi(t)$ is the value, at that $s$, of the function this series calls $\xi$. Volume~5, Theorem~10.8, is the definition:
\[
\xi(s)=\frac12 s(s-1)\pi^{-s/2}\Gamma\Bigl(\frac s2\Bigr)\zeta(s)=\frac12+\frac12 s(s-1)I(s).
\]
A real root of his $\xi(t)$ is a zero of $\zeta$ on the line $\operatorname{Re}s=\frac12$. The epigraph is his sentence about those roots. He adds that he set the search aside after some fleeting vain attempts. Edwards translates the memoir. The German is ``nach einigen fl\"uchtigen vergeblichen Versuchen''.

The $\xi$ of this series is a function of the complex variable $s$. It is entire, and $\xi(1-s)=\xi(s)$. Proposition~\ref{prop:xi} says it has no real zeros. Riemann's sentence is about real $t$, not about real $s$.

In that language the three claims are these. The function $\xi$ is a product over its zeros, as a polynomial is a product over its roots. The number of zeros with ordinate between $0$ and $T$ is about $\frac T{2\pi}\log\frac T{2\pi}-\frac T{2\pi}$. Every one of those zeros lies on the line $\operatorname{Re}s=\frac12$.

A product over no zeros would be a non-zero constant, and the counting formula tends to infinity, so either of the first two claims forces infinitely many zeros. Corollary~\ref{cor:infinitezeros} proves the infinitude directly. The line $\operatorname{Re}s=\frac12$ is the Riemann hypothesis. It is still open.

Hadamard proved the product in 1893, building the theory of entire functions of finite order in order to do it. Three years later he used the product, as de la Vallée Poussin did, to prove the prime number theorem. Von Mangoldt proved the counting formula in 1905.

This chapter proves the two theorems that are proved, and one more. The zeros stay away from the line $\operatorname{Re}s=1$ by a margin of size $c/\log(|t|+2)$. Chapter~\ref{ch:explicit} turns that margin into an error term for the prime number theorem. The tools are three pieces of complex analysis that Volume~5 did not need: Stirling's formula off the real line, Jensen's formula, and the Borel--Carathéodory lemma. The argument principle is Volume~5, Theorem~9.37.

A check of the real formula, before the complex one. At $x=5$, $\Gamma(5)=24$. The three-term expression
\[
\exp\Bigl(\bigl(5-\tfrac12\bigr)\log5-5+\tfrac12\log(2\pi)\Bigr)
\]
equals $23.604$. It is short of $24$ by $0.396$, a relative error of $1.65\%$. Volume~5, Theorem~5.8, is the real case. The same shape for $\sigma+it$ is what the next section proves.

\begin{think}
Figure~\ref{fig:nt} shows $N(T)$, the number of zeros with ordinate in $]0,T]$, tracking Riemann's curve. Before reading on: why would the number of zeros grow like $T\log T$, faster than linearly? Where in the functional equation could a factor $\log T$ come from?
\end{think}

\begin{figure}[ht]
\centering
\fig{ch06-nt}
\caption{The staircase $N(T)$ (charcoal) for $0\le T\le100$, with $29$ steps. These ordinates are sign changes of Hardy's $Z(t)$, the first at $\gamma=14.13$ and the last before $100$ at $\gamma=98.83$. The smooth curve (coral) is Riemann's main term $\frac T{2\pi}\log\frac T{2\pi e}+\frac78$. The difference $S(T)$ stays small: $|S(T)|<0.83$ for $1\le T\le100$.}
\label{fig:nt}
\end{figure}

\section{Stirling's formula in the complex plane}

The functional equation multiplies $\zeta$ by $\pi^{-s/2}\Gamma(s/2)$. A question about $\zeta$ high in the critical strip is, in part, a question about $\Gamma(\sigma+it)$ for large $t$. Volume~5 proved Stirling's formula for real arguments. The complex version is the same formula on a sector. The device is to write down the right holomorphic function and let the identity theorem finish the job.

For $s$ in the slit plane $\C^-=\C\setminus\left]-\infty,0\right]$, let $\log s$ be the principal logarithm. Let $P(u)=\frac12(\{u\}^2-\{u\})$. It is continuous and $1$-periodic, with $-\frac18\le P\le0$, and $P'(u)=\{u\}-\frac12$ away from the integers. Define, for $s\in\C^-$,
\[
\begin{gathered}J(s)=-\int_0^\infty\frac{P(u)}{(u+s)^2}\dd u,\qq F(s)=\Bigl(s-\frac12\Bigr)\log s-s+\frac12\log2\pi+J(s).\end{gathered}
\]

\begin{theorem}[Stirling's formula]\label{thm:stirlingc}
\begin{enumerate}[label=(\alph*)]
\item $F$ is holomorphic on $\C^-$ and $e^{F}=\Gamma$ there. So $F$ is a holomorphic logarithm of $\Gamma$ on $\C^-$, real on $\left]0,\infty\right[$, which we write $\log\Gamma$.
\item Let $0<\delta\le\frac\pi2$. For $|\arg s|\le\pi-\delta$, $|J(s)|\le\frac{C_\delta}{|s|}$ and $|J'(s)|\le\frac{C_\delta}{|s|^2}$, with $C_\delta=(1-\cos\delta)^{-3/2}$.
\item Consequently, for $|\arg s|\le\pi-\delta$ and $|s|\ge1$,
\[
\begin{aligned}\log\Gamma(s)&=\Bigl(s-\frac12\Bigr)\log s-s+\frac12\log2\pi\eqbr+O_\delta\Bigl(\frac1{|s|}\Bigr),\\\frac{\Gamma'}\Gamma(s)&=\log s+O_\delta\Bigl(\frac1{|s|}\Bigr).\end{aligned}
\]
\end{enumerate}
\end{theorem}
\begin{proof}
\emph{Holomorphy.} Let $K\subset\C^-$ be compact, at distance $d>0$ from $\left]-\infty,0\right]$, with $|s|\le S$ on $K$. For $u\ge0$ and $s\in K$, $|u+s|\ge d$. Also $|u+s|\ge u-S\ge u/2$ when $u\ge2S$. So $|u+s|\ge c(1+u)$ for some $c>0$. The integrand of $J$ is continuous in $(u,s)$, holomorphic in $s$, and bounded by $1/(8c^2(1+u)^2)$, which is integrable. Volume~5, Theorem~9.15(b) says that $J$ is holomorphic on $\C^-$. So is $F$.

\emph{The functional equation.} We show $F(s+1)=F(s)+\log s$. Substitute $v=u+1$:
\[
J(s+1)=-\int_1^\infty\frac{P(v)}{(v+s)^2}\dd v.
\]
Integrate by parts, using $P(0)=P(1)=0$:
\[
\begin{aligned}J(s+1)-J(s)&=\int_0^1\frac{P(u)}{(u+s)^2}\dd u\eqbr=\int_0^1\frac{u-\frac12}{u+s}\dd u\\&=1-\Bigl(s+\frac12\Bigr)\eqbr\quad\times\bigl(\log(s+1)-\log s\bigr).\end{aligned}
\]
The algebra uses $\frac{u-1/2}{u+s}=1-\frac{s+1/2}{u+s}$, and $\log(u+s)$ is a primitive of $1/(u+s)$ along the segment $[s,s+1]\subset\C^-$. Add $(s+\frac12)\log(s+1)-(s-\frac12)\log s-1$. The sum is $F(s+1)-F(s)=\log s$.

\emph{Identification.} For real $x>0$, the difference $H(x)=F(x)-\log\Gamma(x)$ is real, and $H(x+1)=H(x)$ by the functional equations of $F$ and of $\Gamma$. As $x\to\infty$,
\[
|J(x)|\le\frac18\int_0^\infty\frac{\mathrm du}{(u+x)^2}=\frac1{8x}\to0,
\]
and $\log\Gamma(x)-(x-\frac12)\log x+x-\frac12\log2\pi\to0$ by Volume~5, Theorem~5.8. The periodic function $H$ tends to $0$, so $H=0$ and $e^F=\Gamma$ on $\left]0,\infty\right[$. Both sides are holomorphic on the connected set $\C^-$ (Volume~5, Theorem~9.36(b) for $\Gamma$). They agree on $\C^-$ by the identity theorem (Volume~5, Theorem~9.18).

(b) Write $s=re^{i\phi}$ with $|\phi|\le\pi-\delta$, and $c=\cos\phi\ge-\cos\delta$. For $u\ge0$,
\[
|u+s|^2=u^2+r^2+2urc\ge(1-\cos\delta)(u^2+r^2).
\]
If $c\ge0$ this is immediate, since $\cos\delta\ge0$. If $c<0$ it follows from $2ur\le u^2+r^2$. With $\kappa=1-\cos\delta$,
\[
\begin{aligned}|J(s)|&\le\frac18\int_0^\infty\frac{\mathrm du}{\kappa(u^2+r^2)}=\frac\pi{16\kappa r},\\|J'(s)|&=\Bigl|2\int_0^\infty\frac{P(u)}{(u+s)^3}\dd u\Bigr|\eqbr\le\frac14\int_0^\infty\frac{\mathrm du}{\kappa^{3/2}(u^2+r^2)^{3/2}}=\frac1{4\kappa^{3/2}r^2}.\end{aligned}
\]
The integral $\int_0^\infty(u^2+r^2)^{-3/2}\mathrm du=1/r^2$ comes from the substitution $u=r\tan\theta$. Both bounds are at most $C_\delta/|s|$ and $C_\delta/|s|^2$.

(c) The first formula is (a) together with (b). Differentiate $F$:
\[
\frac{\Gamma'}{\Gamma}(s)=F'(s)=\log s+\frac{s-1/2}s-1+J'(s)=\log s-\frac1{2s}+J'(s).
\]
\end{proof}

On a vertical line the formula is concrete. Let $\sigma$ be fixed and let $t\to+\infty$. Then
\[
\log(\sigma+it)=\log t+i\frac\pi2-\frac{i\sigma}t+O\Bigl(\frac1{t^2}\Bigr).
\]
Multiply out:
\begin{equation}\label{eq:stirlingvert}
\begin{aligned}\log\Gamma(\sigma+it)&=\Bigl(\sigma-\frac12\Bigr)\log t-\frac{\pi t}2+\frac12\log2\pi\\&\quad+i\Bigl(t\log t-t+\frac\pi2\Bigl(\sigma-\frac12\Bigr)\Bigr)\eqbr+O\Bigl(\frac1t\Bigr).\end{aligned}
\end{equation}
So $|\Gamma(\sigma+it)|\sim\sqrt{2\pi}\,t^{\sigma-1/2}e^{-\pi t/2}$. Volume~5, Theorem~9.36(d), is the same exponential decay on the line $\operatorname{Re}z=\frac12$. The formula above puts it on every vertical line, together with the power of $t$ that decides the rest of the chapter.

One numerical check of \eqref{eq:stirlingvert}. At $s=\frac14+50i$, $|\Gamma(s)|=7.327\times10^{-35}$. The main term $\sqrt{2\pi}\,t^{\sigma-1/2}e^{-\pi t/2}$ agrees with this modulus to a relative error of $3\times10^{-6}$.

\section{The function \texorpdfstring{$\xi$}{xi}}

Riemann's $\xi$ packs $\zeta$, the Gamma factor and the pole into one entire function. Recall it from Volume~5, Theorem~10.8:
\[
\xi(s)=\frac12 s(s-1)\pi^{-s/2}\Gamma\Bigl(\frac s2\Bigr)\zeta(s).
\]
That theorem shows $\xi(s)=\frac12+\frac12 s(s-1)I(s)$, where
\[
I(s)=\int_1^\infty\omega(t)\bigl(t^{s/2}+t^{(1-s)/2}\bigr)\frac{\mathrm dt}t
\]
is entire, and that $\xi(1-s)=\xi(s)$. Jacobi's theta, which produces $\omega$, is $\vartheta$. The symbol $\Lambda$ is von Mangoldt's function.

\begin{proposition}\label{prop:xi}
\begin{enumerate}[label=(\alph*)]
\item $\xi$ is entire, $\xi(s)=\xi(1-s)$, $\xi(\overline s)=\overline{\xi(s)}$, and $\xi(0)=\xi(1)=\frac12$.
\item The zeros of $\xi$ are exactly the zeros $\rho=\beta+i\gamma$ of $\zeta$ with $0<\beta<1$, with the same multiplicities. They are symmetric with respect to the real axis and to the line $\operatorname{Re}s=\frac12$, and none is real.
\item There is a constant $C$ such that $|\xi(s)|\le e^{C|s|\log|s|}$ for $|s|\ge2$.
\end{enumerate}
\end{proposition}
\begin{proof}
(a) The formula $\xi(s)=\frac12+\frac12 s(s-1)I(s)$ of Volume~5, Theorem~10.8, shows that $\xi$ is entire, symmetric under $s\mapsto1-s$, and equal to $\frac12$ at $0$ and at $1$. Conjugation symmetry holds for $\operatorname{Re}s>1$, where every factor is real on the real axis, and therefore everywhere, by the identity theorem.

(b) For $\operatorname{Re}s>1$, the Euler product gives $\zeta(s)\ne0$, and $\Gamma(s/2)\ne0$ (Volume~5, Theorem~9.36(c)), so $\xi(s)\ne0$. By the symmetry $\xi(s)=\xi(1-s)$, one has $\xi\ne0$ on $\operatorname{Re}s<0$. On $\operatorname{Re}s=1$, Theorem~\ref{thm:line} gives $\zeta(1+it)\ne0$ for $t\ne0$, and $\xi(1)=\frac12$. The symmetry removes the line $\operatorname{Re}s=0$ as well. In the open strip, $s(s-1)\pi^{-s/2}\Gamma(s/2)$ is holomorphic and non-zero, so $\xi$ and $\zeta$ have the same zeros there, with the same orders. Finally $\zeta(\sigma)<0$ for $0<\sigma<1$ (Volume~5, Exercise~10.3), so no zero is real.

(c) If $\sigma<\frac12$, then $|s|\le|1-s|\le|s|+1$. By symmetry it is enough to treat $\sigma\ge\frac12$ and $|s|\ge2$, where $|s-1|\ge1$. The proof of Volume~5, Proposition~10.6, gives, for $\sigma>0$,
\[
\Bigl|\zeta(s)-\frac1{s-1}\Bigr|\le|s|\sum_n n^{-\sigma-1}\le|s|\Bigl(1+\frac1\sigma\Bigr),
\]
so $|\zeta(s)|\le1+3|s|$ when $\sigma\ge\frac12$ and $|s-1|\ge1$. Next $|\pi^{-s/2}|\le1$, and $|\Gamma(s/2)|\le\Gamma(\sigma/2)$ (Volume~5, Theorem~9.36(a)). For $x\ge1$ and every $t>0$, $t^{x-1}e^{-t/2}\le(2x)^{x-1}$, because the maximum is at $t=2(x-1)$. So
\[
\Gamma(x)\le(2x)^{x-1}\int_0^\infty e^{-t/2}\dd t\le(2x)^x.
\]
For $\frac14\le x\le1$, $\Gamma(x)=\Gamma(x+1)/x\le1/x\le4$, because $\Gamma\le1$ on $[1,2]$ (it is convex there, with $\Gamma(1)=\Gamma(2)=1$). Altogether, on the region under consideration,
\[
|\xi(s)|\le|s|^2(1+3|s|)\bigl(4+(2|s|)^{|s|}\bigr),
\]
which is at most $e^{C|s|\log|s|}$ for $|s|\ge2$.
\end{proof}

The bound (c) is sharp in shape. For real $\sigma\to\infty$,
\[
\xi(\sigma)\ge\frac12\sigma(\sigma-1)\pi^{-\sigma/2}\Gamma\Bigl(\frac\sigma2\Bigr),
\]
which grows like $e^{\frac12\sigma\log\sigma\,(1+o(1))}$ by Stirling. We say that $\xi$ has \emph{order $1$}. That single number makes the next three sections work.

Figure~\ref{fig:zfree} draws the symmetries in (b), and the margin proved at the end of the chapter. The four dots are one zero and the three partners forced by $\xi(s)=\xi(1-s)$ and by conjugation. They are a schematic, not a measured ordinate. The ordinates that have been computed are the ones in Figure~\ref{fig:nt}.

\begin{figure}[ht]
\centering
\fig{ch06-new-strip}
\caption{The critical strip $0\le\sigma\le1$ (charcoal tint) and the half-plane $\sigma\ge1$ (coral tint), where Theorem~\ref{thm:line} already forbids zeros. The coral curve is the boundary $\sigma=1-c/\log(|t|+2)$, drawn schematically: the gap from the line $\sigma=1$ shrinks as $|t|$ grows. A zero $\rho$ off the critical line brings three partners with it.}
\label{fig:zfree}
\end{figure}

\section{Counting zeros with a circle}

How can one count the zeros of a function without finding them? In 1899 Jensen, an engineer at the Copenhagen telephone company who did mathematics in his spare time, gave the answer: average $\log|f|$ over a circle. Figure~\ref{fig:jensencirc} is the picture. Each zero at distance $|z_k|$ from the centre contributes $\log(R/|z_k|)$ to that average.

\begin{figure}[ht]
\centering
\fig{ch06-new-jensen}
\caption{Jensen's formula compares $\log|f(0)|$ with the average of $\log|f|$ on the circle $|z|=R$. Each zero inside the circle raises the average by $\log(R/|z_k|)$. A function that is not allowed to grow fast is not allowed many zeros.}
\label{fig:jensencirc}
\end{figure}

\begin{theorem}[Jensen's formula]\label{thm:jensen}
Let $f$ be holomorphic on an open set containing the closed disc $|z|\le R$, with $f(0)\ne0$, and let $z_1,\dots,z_n$ be its zeros in $|z|<R$, repeated according to multiplicity. If $f$ has no zero on $|z|=R$, then
\[
\begin{aligned}\log|f(0)|&+\sum_{k=1}^n\log\frac R{|z_k|}\eqbr=\frac1{2\pi}\int_0^{2\pi}\log|f(Re^{i\theta})|\dd\theta.\end{aligned}
\]
\end{theorem}
\begin{proof}
For $|a|<R$, the Blaschke factor $B_a(z)=\frac{R(z-a)}{R^2-\overline a z}$ is holomorphic on the disc $|z|<R^2/|a|$, which is larger than $|z|\le R$. It vanishes only at $a$, and it has modulus $1$ on $|z|=R$. There $R^2=z\bar z$, so $|R^2-\overline a z|=|z|\,|\bar z-\bar a|=R|z-a|$. Let $g=f/\prod_k B_{z_k}$. The zeros cancel, with their multiplicities, so $g$ is holomorphic on an open disc $D$ slightly larger than $|z|\le R$. It has no zeros on the closed disc, hence none on a slightly smaller open disc $D'\supset\{|z|\le R\}$. On the convex set $D'$, the quotient $g'/g$ has a primitive $\ell$ (Volume~5, Theorem~9.9). Then $(ge^{-\ell})'=0$, so after adding a constant one has $e^\ell=g$. Thus $\ell$ is a holomorphic logarithm of $g$. The mean value property, which is Cauchy's formula at the centre, gives
\[
\ell(0)=\frac1{2\pi}\int_0^{2\pi}\ell(Re^{i\theta})\dd\theta.
\]
Take real parts: $\log|g(0)|$ equals the average of $\log|g(Re^{i\theta})|$. On the circle $|g|=|f|$, and $|g(0)|=|f(0)|\prod_k R/|z_k|$ because $|B_a(0)|=|a|/R$.
\end{proof}

We use the formula as an inequality. The inequality needs no condition on the circle itself.

\begin{corollary}\label{cor:jensen}
Let $f$ be holomorphic on an open set containing $|z|\le R$, with $f(0)\ne0$ and $|f|\le M$ on $|z|=R$. For $0<r<R$, the number $n(r)$ of zeros of $f$ in $|z|\le r$ satisfies
\[
n(r)\log\frac Rr\le\log\frac M{|f(0)|}.
\]
\end{corollary}
\begin{proof}
The zeros in the compact disc are finite in number, so there are radii $R'\in\left]r,R\right[$, arbitrarily close to $R$, with no zero on $|z|=R'$. In Jensen's formula for $R'$, each of the $n(r)$ zeros in $|z|\le r$ contributes at least $\log(R'/r)$. The other zeros contribute non-negative terms. The right side is at most $\log\max_{|z|=R'}|f|\le\log M$, by the maximum principle (Volume~5, Theorem~9.19). Let $R'\to R$.
\end{proof}

\begin{beauty}
Jensen's formula is a conservation law. The average of $\log|f|$ on a circle exceeds its value at the centre by exactly the amount that the zeros inside cost. A zero at distance $|z_k|$ costs $\log(R/|z_k|)$. A function that is not allowed to grow fast is not allowed many zeros.
\end{beauty}

Apply this to $\xi$, with $|\xi(0)|=\frac12$ and Proposition~\ref{prop:xi}(c) on the circle of radius $2r$. Corollary~\ref{cor:jensen} gives
\begin{equation}\label{eq:nr}
\#\{\rho:|\rho|\le r\}\le\frac{\log(2e^{2Cr\log2r})}{\log2}=O(r\log r).
\end{equation}
Summing by parts, $\sum_\rho|\rho|^{-1-\varepsilon}<\infty$ for every $\varepsilon>0$. In particular $\sum_\rho|\rho|^{-2}<\infty$. Let $r_0>0$ be the smallest modulus of a zero. It exists, because $\xi(0)\ne0$ and the zeros are isolated. Then
\[
\sum_{|\rho|\le R}|\rho|^{-2}=\int_{r_0^-}^R\frac{\mathrm dn(r)}{r^2}=\frac{n(R)}{R^2}+2\int_{r_0}^R\frac{n(r)}{r^3}\dd r,
\]
which stays bounded as $R\to\infty$. Exercise~\ref{ex:zerosums} shows that $\sum|\rho|^{-1}$ diverges. There are many zeros, and not too many.

\section{Hadamard's product}

A polynomial is determined, up to a constant, by its roots. An entire function is not: $e^z$ has no zeros at all. Hadamard's discovery is that if the function does not grow too fast, the ambiguity is only an exponential factor of controlled degree. We prove this for $\xi$, in the form of a partial fraction expansion of $\xi'/\xi$. The product follows. The tool is a lemma that bounds a holomorphic function by its real part alone.

\begin{lemma}[Borel--Carathéodory]\label{lem:borel}
Let $h$ be holomorphic on an open set containing $|z|\le R$, with $h(0)=0$ and $\operatorname{Re}h\le A$ on $|z|=R$. Then for $|z|\le r<R$ and $k\ge1$,
\[
|h^{(k)}(z)|\le\frac{2A\,k!\,R}{(R-r)^{k+1}}.
\]
\end{lemma}
\begin{proof}
The function $h$ is holomorphic on a disc $|z|<R_1$ with $R_1>R$, so $h(z)=\sum_{n\ge1}a_nz^n$ there (Volume~5, Theorem~9.11), and the series converges uniformly on $|z|=R$. Integrate term by term against $e^{-in\theta}$ on the circle:
\[
\begin{gathered}a_nR^n=\frac1{2\pi}\int_0^{2\pi}h(Re^{i\theta})e^{-in\theta}\dd\theta,\qq0=\frac1{2\pi}\int_0^{2\pi}\overline{h(Re^{i\theta})}e^{-in\theta}\dd\theta\end{gathered}
\]
for $n\ge1$. The second integral vanishes because its conjugate is the coefficient of $z^{-n}$, which is $0$. Adding,
\[
a_nR^n=\frac1\pi\int_0^{2\pi}\operatorname{Re}h(Re^{i\theta})e^{-in\theta}\dd\theta.
\]
Since $\int_0^{2\pi}e^{-in\theta}\dd\theta=0$, one also has
\[
a_nR^n=-\frac1\pi\int_0^{2\pi}\bigl(A-\operatorname{Re}h(Re^{i\theta})\bigr)e^{-in\theta}\dd\theta.
\]
Now $A-\operatorname{Re}h\ge0$ on the circle, and its integral is $2\pi A$, because $\int\operatorname{Re}h=2\pi\operatorname{Re}h(0)=0$. So $|a_n|R^n\le2A$, and
\[
\begin{aligned}|h^{(k)}(z)|&\le\sum_{n\ge k}n(n-1)\cdots(n-k+1)\frac{2A}{R^n}r^{n-k}\eqbr=\frac{2A}{R^k}\cdot\frac{k!}{(1-r/R)^{k+1}},\end{aligned}
\]
using $\sum_{n\ge k}\binom nkx^{n-k}=(1-x)^{-k-1}$. This is the claimed bound.
\end{proof}

Only an upper bound on $\operatorname{Re}h$ is used. For $h=\log(f/f(0))$ that upper bound is an upper bound on $|f|$, which is what a growth estimate gives.

\begin{lemma}\label{lem:localzeros}
Let $f$ be holomorphic on an open set containing $|z|\le R$, with $f(0)\ne0$ and $|f(z)|\le M|f(0)|$ for $|z|\le R$, where $M\ge2$. Let $\mathcal Z$ be the set of zeros of $f$ in $|z|\le R/2$, with multiplicity. Then $\#\mathcal Z\le\frac{\log M}{\log2}$, and for $|z|\le R/4$, $z\notin\mathcal Z$,
\[
\begin{gathered}\Bigl|\frac{f'}f(z)-\sum_{\rho\in\mathcal Z}\frac1{z-\rho}\Bigr|\le\frac{17\log M}R,\\\Bigl|\Bigl(\frac{f'}f\Bigr)'(z)+\sum_{\rho\in\mathcal Z}\frac1{(z-\rho)^2}\Bigr|\le\frac{129\log M}{R^2}.\end{gathered}
\]
\end{lemma}
\begin{proof}
The count is Corollary~\ref{cor:jensen} with $r=R/2$. Let
\[
g(z)=f(z)\prod_{\rho\in\mathcal Z}\frac{R^2-\overline\rho z}{R(z-\rho)}.
\]
Each factor has modulus $1$ on $|z|=R$, as the Blaschke factors do, and its pole at $\rho$ is cancelled by a zero of $f$. Its zero $R^2/\overline\rho$ has modulus at least $2R$. So $g$ is holomorphic near $|z|\le R$, and $|g|=|f|\le M|f(0)|$ on $|z|=R$, hence on the disc by the maximum principle. Also $|g(0)|=|f(0)|\prod R/|\rho|\ge|f(0)|$. The function $g$ has no zero in $|z|\le R/2$. As in the proof of Jensen's formula, $h=\log(g/g(0))$ is holomorphic near $|z|\le R/2$, with $h(0)=0$ and $\operatorname{Re}h=\log(|g|/|g(0)|)\le\log M$. Lemma~\ref{lem:borel}, with radius $R/2$ and with $r=R/4$, gives $|h'|\le16\log M/R$ and $|h''|\le128\log M/R^2$ on $|z|\le R/4$. Now
\[
h'=\frac{f'}f-\sum_{\rho\in\mathcal Z}\frac1{z-\rho}-\sum_{\rho\in\mathcal Z}\frac{\overline\rho}{R^2-\overline\rho z}.
\]
For $|z|\le R/4$ and $|\rho|\le R/2$,
\[
\Bigl|\frac{\overline\rho}{R^2-\overline\rho z}\Bigr|\le\frac{R/2}{R^2-R^2/8}=\frac4{7R}.
\]
There are at most $(\log M)/\log2$ such terms, so the last sum is at most
\[
\frac{4}{7\log2}\frac{\log M}R<\frac{\log M}R.
\]
Differentiating once more, the corresponding terms are at most $16/(49R^2)$ each, and $16/(49\log2)<1$.
\end{proof}

\begin{theorem}[Hadamard, 1893]\label{thm:hadamard}
Let $\rho$ run through the zeros of $\xi$, with multiplicity, and let $B=\xi'(0)/\xi(0)$. Then, for every $s$ that is not a zero,
\[
\frac{\xi'(s)}{\xi(s)}=B+\sum_\rho\Bigl(\frac1{s-\rho}+\frac1\rho\Bigr),
\]
the series converging absolutely and uniformly on compact sets avoiding the zeros. Equivalently,
\[
\xi(s)=\frac12 e^{Bs}\prod_\rho\Bigl(1-\frac s\rho\Bigr)e^{s/\rho}.
\]
\end{theorem}
\begin{proof}
\emph{The second derivative.} Fix $s$, not a zero, and let $R\ge4|s|+8$. On $|z|\le R$, Proposition~\ref{prop:xi}(c) gives $|\xi(z)|\le e^{CR\log R}=2e^{CR\log R}|\xi(0)|$, so Lemma~\ref{lem:localzeros} applies with $\log M=O(R\log R)$:
\[
\Bigl(\frac{\xi'}{\xi}\Bigr)'(s)=-\sum_{|\rho|\le R/2}\frac1{(s-\rho)^2}+O\Bigl(\frac{\log R}R\Bigr).
\]
As $R\to\infty$ the error tends to $0$, and the sum tends to the full series. The series converges absolutely: for $|\rho|\ge2|s|$ one has $|s-\rho|\ge|\rho|/2$, and $\sum|\rho|^{-2}<\infty$. So $(\xi'/\xi)'=-\sum_\rho(s-\rho)^{-2}$.

\emph{Integrating once.} Let
\[
\Phi(s)=\sum_\rho\Bigl(\frac1{s-\rho}+\frac1\rho\Bigr)=\sum_\rho\frac s{\rho(s-\rho)}.
\]
On a compact set $K$ that avoids the zeros, with $|s|\le S$, the terms with $|\rho|\ge2S$ are at most $2S/|\rho|^2$. The series converges uniformly on $K$, so $\Phi$ is holomorphic off the zeros, and $\Phi'=-\sum(s-\rho)^{-2}$ (Volume~5, Theorem~9.15(a)). Hence $\xi'/\xi-\Phi$ has derivative $0$ on the complement of the zeros. That complement is connected, because the zeros are isolated, so the difference is a constant. Its value at $s=0$ is $B$.

\emph{The product.} For $|w|\le\frac12$, the principal value $\ell(w)=\log(1-w)+w=-\sum_{k\ge2}w^k/k$ satisfies $|\ell(w)|\le|w|^2$. For $|s|\le S$ and $|\rho|\ge2S$, the factor $(1-s/\rho)e^{s/\rho}$ is $e^{\ell(s/\rho)}$, and
\[
\sum|\ell(s/\rho)|\le S^2\sum|\rho|^{-2}.
\]
The product converges uniformly on $|s|\le S$ to a holomorphic function $E(s)$, whose zeros are the $\rho$. Differentiating the uniformly convergent sum of logarithms, and the finitely many remaining factors directly, gives $E'/E=\Phi$. Then $\xi/E$ is entire and has no zeros, and its logarithmic derivative is $B$, so $\xi=\xi(0)e^{Bs}E$.
\end{proof}

\begin{corollary}\label{cor:infinitezeros}
$\zeta$ has infinitely many zeros in the strip $0<\operatorname{Re}s<1$.
\end{corollary}
\begin{proof}
If there were finitely many, then $E$ would be a polynomial times $e^{cs}$, and $|\xi(\sigma)|\le e^{C'\sigma}$ for real $\sigma\ge1$. But
\[
\xi(\sigma)\ge\frac12\sigma(\sigma-1)\pi^{-\sigma/2}\Gamma\Bigl(\frac\sigma2\Bigr)
\]
grows faster than any exponential, by Stirling's formula.
\end{proof}

The logarithmic derivative of the definition of $\xi$ is the form in which the theorem is used. With $\psi_\Gamma=\Gamma'/\Gamma$ (the letter $\psi$ is already von Mangoldt's function),
\begin{equation}\label{eq:zetapf}
\begin{aligned}\frac{\zeta'}{\zeta}(s)&=B+\sum_\rho\Bigl(\frac1{s-\rho}+\frac1\rho\Bigr)-\frac1s-\frac1{s-1}\eqbr+\frac12\log\pi-\frac12\psi_\Gamma\Bigl(\frac s2\Bigr).\end{aligned}
\end{equation}
The pole of $\zeta$ appears as $-1/(s-1)$. The trivial zeros come from the poles of $\psi_\Gamma(s/2)$ at $0,-2,-4,\dots$. The pole at $0$ cancels $-1/s$. Every non-trivial zero contributes a simple pole. Exercises~\ref{ex:digammahalf} and~\ref{ex:bsum} show that
\[
B=\frac12\log4\pi-1-\frac\gamma2=-0.0230957\ldots=-\sum_\rho\operatorname{Re}\frac1\rho.
\]

\section{The Riemann--von Mangoldt formula}

Order the ordinates $\gamma>0$ of the zeros, and let $N(T)$ be the number of zeros $\rho=\beta+i\gamma$ with $0<\gamma\le T$, counted with multiplicity. Following Volume~5, Chapter~10, let
\[
\alpha(T)=\operatorname{Im}\log\Gamma\Bigl(\frac14+\frac{iT}2\Bigr)-\frac T2\log\pi,
\]
with the logarithm of Theorem~\ref{thm:stirlingc}. This is the phase $\alpha$ of Hardy's function in Volume~5. The same phase is often written $\theta$ in other books. This series keeps $\theta$ for Chebyshev's function, and it keeps $\alpha$ for the phase. By \eqref{eq:stirlingvert} with $\sigma=\frac14$ and $t=T/2$,
\begin{equation}\label{eq:theta}
\alpha(T)=\frac T2\log\frac T{2\pi}-\frac T2-\frac\pi8+O\Bigl(\frac1T\Bigr).
\end{equation}

\begin{theorem}[Riemann 1859, von Mangoldt 1905]\label{thm:rvm}
For $T\ge3$ not the ordinate of a zero,
\[
\begin{gathered}N(T)=\frac{\alpha(T)}\pi+1+S(T),\qq S(T)=\frac1\pi\operatorname{Im}\int_L\frac{\zeta'(s)}{\zeta(s)}\dd s,\end{gathered}
\]
where $L$ is the path from $2$ to $2+iT$ to $\frac12+iT$, and $S(T)=O(\log T)$. Hence
\[
N(T)=\frac T{2\pi}\log\frac T{2\pi e}+\frac78+O(\log T).
\]
\end{theorem}

The constant $\frac78$ is part of the formula. Riemann's main term receives $\frac78$ from the Gamma factor and the pole. The content of the theorem is the error $O(\log T)$, set against a main term that grows like $T\log T$.

At $T=100$ the main term equals $29.002$. The staircase in Figure~\ref{fig:nt} has $29$ steps. The $O(\log T)$ term is small on this range: the caption records $|S(T)|<0.83$ up to $100$.

\begin{proof}
\emph{Counting with the argument principle.} Let $D$ be the open rectangle $-1<\sigma<2$, $0<t<T$. The zeros of $\xi$ in $D$ are the zeros of $\zeta$ with $0<\gamma<T$. None is real, and none lies on $\partial D$. By Volume~5, Theorem~9.37,
\[
N(T)=\frac1{2\pi i}\int_{\partial D}\frac{\xi'}{\xi}.
\]
The integral is real, so $2\pi N(T)=\operatorname{Im}\int_{\partial D}\xi'/\xi$. On the bottom side, $\xi$ is real and positive: $\xi(0)=\frac12$, and there are no real zeros, so $\operatorname{Im}(\xi'/\xi)=0$ there. The map $s\mapsto1-\overline s$ exchanges the right half of the rest of $\partial D$, which is the path $L$, with the left half, and it reverses direction. Since $\xi(1-\overline s)=\overline{\xi(s)}$, the imaginary part of $\xi'/\xi$ integrates to the same value on both halves. So $\pi N(T)=\operatorname{Im}\int_L\xi'/\xi$.

\emph{The four factors.} The logarithmic derivative splits as
\[
\frac{\xi'}{\xi}=\frac1s+\frac1{s-1}-\frac12\log\pi+\frac12\psi_\Gamma\Bigl(\frac s2\Bigr)+\frac{\zeta'}{\zeta}.
\]
Along $L$:
\begin{itemize}
\item $\operatorname{Im}\int_L\bigl(\frac1s+\frac1{s-1}\bigr)$ is the change in a continuous argument of $s(s-1)$. The product starts at $2>0$ and ends at
\[
\bigl(\tfrac12+iT\bigr)\bigl(-\tfrac12+iT\bigr)=-T^2-\tfrac14<0.
\]
In between, $\operatorname{Im}s(s-1)=t(2\sigma-1)>0$. The change is exactly $\pi$.
\item $\operatorname{Im}\int_L(-\frac12\log\pi)\dd s=-\frac12\log\pi\cdot\operatorname{Im}(\frac12+iT-2)=-\frac T2\log\pi$.
\item $\int_L\frac12\psi_\Gamma(s/2)\dd s=\log\Gamma(\frac14+iT/2)-\log\Gamma(1)$. The imaginary part is $\operatorname{Im}\log\Gamma(\frac14+iT/2)$, because $L$ stays in $\C^-$.
\end{itemize}
Adding these contributions, $\pi N(T)=\pi+\alpha(T)+\pi S(T)$.

\emph{The bound on $S(T)$.} On the vertical side,
\[
\operatorname{Re}\zeta(2+it)\ge1-\sum_{n\ge2}n^{-2}=2-\zeta(2)>0.35,
\]
so a continuous argument of $\zeta$ stays in $\left]-\pi/2,\pi/2\right[$ and changes by less than $\pi$. On the horizontal side, suppose $\operatorname{Re}\zeta(\sigma+iT)$ vanishes at $q$ points of $[\frac12,2]$. Those points cut the segment into at most $q+1$ pieces. On each piece $\operatorname{Re}\zeta$ has constant sign, so the argument changes by less than $\pi$. Thus $\pi|S(T)|<\pi+(q+1)\pi$.

To bound $q$, let $k(z)=\frac12\bigl(\zeta(z+iT)+\zeta(z-iT)\bigr)$. It is holomorphic on $|z-2|\le\frac74$, because the poles of the two terms are at $1\mp iT$, far from that disc. For real $z=\sigma$, one has $\zeta(\sigma-iT)=\overline{\zeta(\sigma+iT)}$, so $k(\sigma)=\operatorname{Re}\zeta(\sigma+iT)$. The integer $q$ is at most the number of zeros of $k$ in $|z-2|\le\frac32$. On $|z-2|\le\frac74$, one has $\operatorname{Re}(z\pm iT)\ge\frac14$ and $|z\pm iT|\le T+4$, so $|k(z)|\le1+5(T+4)$ by the bound in the proof of Proposition~\ref{prop:xi}(c), which gives $|\zeta(s)|\le1/|s-1|+5|s|$ for $\sigma\ge\frac14$. And $k(2)>0.35$. Corollary~\ref{cor:jensen} gives
\[
q\le\frac{\log\bigl(3(5T+21)\bigr)}{\log(7/6)}=O(\log T).
\]
Finally, \eqref{eq:theta} turns $\alpha(T)/\pi+1$ into
\[
\frac T{2\pi}\log\frac T{2\pi}-\frac T{2\pi}-\frac18+1+O\Bigl(\frac1T\Bigr).
\]
\end{proof}

\begin{corollary}\label{cor:nunit}
For $T\ge3$, the number of zeros with $T<\gamma\le T+1$ is $O(\log T)$. For every $s=\sigma+it$ with $|t|\ge2$,
\[
\sum_\rho\frac1{1+(t-\gamma)^2}=O(\log|t|).
\]
\end{corollary}
\begin{proof}
The difference is $N(T+1)-N(T)=(\alpha(T+1)-\alpha(T))/\pi+O(\log T)$. Theorem~\ref{thm:stirlingc}(c), differentiated under the imaginary part, gives $\alpha'(u)=\frac12\log(u/(2\pi))+O(1/u)$. For the sum, the zeros are symmetric about the real axis, so the sum is at most twice the sum over $\gamma>0$. Group the ordinates with $k\le|\gamma-t|<k+1$, for $k\ge0$. There are $O(\log(|t|+k+2))$ of them, contributing $O\bigl(\log(|t|+k+2)/(1+k^2)\bigr)$, and
\[
\sum_k\frac{\log(|t|+k+2)}{1+k^2}=O(\log|t|).
\]
\end{proof}

\begin{remark}
On the critical line, $\xi(\frac12+it)$ is real, and its sign changes are zeros on the line. Volume~5 found them through Hardy's function $Z(t)=e^{i\alpha(t)}\zeta(\frac12+it)$. Theorem~\ref{thm:rvm} counts every zero in the strip. If a computation finds $N(T)$ sign changes of $Z$ on $]0,T]$, then every zero up to height $T$ lies on the line and is simple. Making the comparison rigorous needs care with $S(T)$ at the endpoint. Turing found the method in 1953, using Littlewood's bound $\int_0^TS=O(\log T)$, quoted from Titchmarsh, Chapter~IX. By this method Gourdon checked the first $10^{13}$ zeros in 2004. For a first taste: $Z$ changes sign $79$ times on $]0,200]$, and Theorem~\ref{thm:rvm}, with $S$ computed to enough accuracy, gives $N(200)=79$. The ordinates used in Figure~\ref{fig:nt} are of this kind. This chapter does not recompute them.
\end{remark}

\section{Zeros near a point, and away from the line}

The partial fraction \eqref{eq:zetapf} equates two expressions that each look divergent. A difference makes the identity local. Near a point $s$, the quotient $\zeta'/\zeta$ is the sum of $1/(s-\rho)$ over the zeros close to $s$, up to an error $O(\log|t|)$.

\begin{proposition}\label{prop:nearzeros}
Uniformly for $-1\le\sigma\le2$ and $|t|\ge2$,
\[
\frac{\zeta'}{\zeta}(\sigma+it)=\sum_{|\gamma-t|\le1}\frac1{s-\rho}+O(\log|t|).
\]
\end{proposition}
\begin{proof}
Let $s_1=2+it$, and subtract \eqref{eq:zetapf} at $s_1$ from \eqref{eq:zetapf} at $s$. The left side at $s_1$ is bounded, by $\sum\Lambda(n)n^{-2}$. The differences $1/s-1/s_1$ and $1/(s-1)-1/(s_1-1)$ are bounded, since $|t|\ge2$. Both $s/2$ and $s_1/2$ lie in the sector $|\arg z|\le\pi-\arctan2$, with modulus at least $1$. Theorem~\ref{thm:stirlingc}(c) gives
\[
\psi_\Gamma\Bigl(\frac s2\Bigr)-\psi_\Gamma\Bigl(\frac{s_1}2\Bigr)=\log\frac s{s_1}+O(1)=O(1).
\]
What remains is $\sum_\rho\bigl(1/(s-\rho)-1/(s_1-\rho)\bigr)$. If $|\gamma-t|>1$, the general term is
\[
\frac{s_1-s}{(s-\rho)(s_1-\rho)},
\]
of modulus at most $3/|\gamma-t|^2\le6/(1+(\gamma-t)^2)$. Corollary~\ref{cor:nunit} says these add up to $O(\log|t|)$. If $|\gamma-t|\le1$, then $|s_1-\rho|\ge2-\beta>1$, so $|1/(s_1-\rho)|<1$, and there are $O(\log|t|)$ such $\rho$. What remains is $\sum_{|\gamma-t|\le1}1/(s-\rho)$.
\end{proof}

Chapter~\ref{ch:pnt} showed that no zero lies on $\operatorname{Re}s=1$. A zero could still approach the line as $t$ grows. The next theorem says how fast it may approach. The constant in the proof is explicit: for large $|t|$ the gap is at least $1/(14A'\log|t|)$, where $A'\ge1$ is the constant collected from the estimates.

\begin{theorem}[de la Vallée Poussin, 1899]\label{thm:zerofree}
There is an absolute constant $c>0$ such that $\zeta(\sigma+it)\ne0$ for
\[
\sigma\ge1-\frac c{\log(|t|+2)}.
\]
\end{theorem}
\begin{proof}
\emph{Three inequalities for $1<\sigma\le2$.} First, the inequality $3+4\cos\phi+\cos2\phi\ge0$ of Theorem~\ref{thm:line}, applied with $\phi=t\log n$ and summed against $\Lambda(n)n^{-\sigma}$ (Lemma~\ref{lem:logderiv}), gives
\[
\begin{aligned}3\Bigl(-\frac{\zeta'}{\zeta}(\sigma)\Bigr)&+4\operatorname{Re}\Bigl(-\frac{\zeta'}{\zeta}(\sigma+it)\Bigr)\eqbr+\operatorname{Re}\Bigl(-\frac{\zeta'}{\zeta}(\sigma+2it)\Bigr)\ge0.\end{aligned}
\]
Second, $-\zeta'/\zeta(s)-1/(s-1)$ is holomorphic near $1$ (the proof of Theorem~\ref{thm:pnt}), hence bounded on $[1,2]$: $-\zeta'/\zeta(\sigma)\le1/(\sigma-1)+A$. Third, for $|t|\ge2$, Proposition~\ref{prop:nearzeros} gives
\[
\operatorname{Re}\Bigl(-\frac{\zeta'}{\zeta}(\sigma+it)\Bigr)=-\sum_{|\gamma-t|\le1}\operatorname{Re}\frac1{s-\rho}+O(\log|t|).
\]
Every term of the sum is positive, because $\operatorname{Re}1/(s-\rho)=(\sigma-\beta)/|s-\rho|^2$ with $\sigma>1>\beta$. Dropping positive terms can only raise an upper bound. So $\operatorname{Re}(-\zeta'/\zeta(\sigma+2it))\le A\log|t|$, and if $\rho_0=\beta_0+it$ is a zero,
\[
\operatorname{Re}\Bigl(-\frac{\zeta'}{\zeta}(\sigma+it)\Bigr)\le-\frac1{\sigma-\beta_0}+A\log|t|,
\]
for a suitable constant $A\ge1$. Collect every constant into one $A'\ge1$, using $\log|t|\ge\log2$.

\emph{The squeeze.} For a zero $\rho_0=\beta_0+it$ with $|t|\ge2$, and for every $\sigma\in\left]1,2\right]$,
\[
\frac4{\sigma-\beta_0}\le\frac3{\sigma-1}+A'\log|t|.
\]
This is the $3$-$4$-$1$ argument of Chapter~\ref{ch:pnt}, made quantitative. The pole contributes $3$. A zero near $1+it$ would contribute $-4$, and $4>3$. Take $\sigma=1+\delta/\log|t|$ with $\delta=1/(2A')$, so that $\sigma\le2$. Then
\[
\frac4{\sigma-\beta_0}\le\Bigl(\frac3\delta+A'\Bigr)\log|t|=7A'\log|t|,
\]
hence $\sigma-\beta_0\ge4/(7A'\log|t|)=8\delta/(7\log|t|)$, and
\[
\begin{aligned}1-\beta_0&=(\sigma-\beta_0)-(\sigma-1)\eqbr\ge\frac{\delta}{7\log|t|}=\frac1{14A'\log|t|}.\end{aligned}
\]
For $|t|<2$ there are finitely many zeros with $0<\beta<1$. They are isolated in a compact set, and none lies on $\operatorname{Re}s=1$, so each satisfies $\beta\le\beta_*<1$. A smaller positive $c$ covers them, and it covers the large-$|t|$ bound as well, because $\log(|t|+2)$ is comparable to $\log|t|$ for $|t|\ge2$.
\end{proof}

\begin{remark}
The region has been widened only slightly in more than a century. Vinogradov and Korobov (1958) reached
\[
\sigma\ge1-\frac c{(\log|t|)^{2/3}(\log\log|t|)^{1/3}},
\]
which gives the best known error term in the prime number theorem (Chapter~\ref{ch:explicit}). The shape is quoted from Titchmarsh, Chapter~VI. The Riemann hypothesis would give the half-plane $\sigma>\frac12$.

Hardy proved in 1914 that infinitely many zeros lie on the line. Selberg proved in 1942 that a positive proportion of them do. Later proportions are quotations, not theorems of this chapter: more than one third (Levinson, 1974), more than two fifths (Conrey, 1989), more than five twelfths (Pratt, Robles, Zaharescu and Zeindler, 2020). The sources are at the end of the chapter.
\end{remark}

\section{What we proved, and what comes next}

Two of Riemann's three claims are theorems. The zeros of $\xi$ are the non-trivial zeros of $\zeta$, and $\xi$ is a product over them (Theorem~\ref{thm:hadamard}): its logarithmic derivative is $B+\sum_\rho\bigl(1/(s-\rho)+1/\rho\bigr)$. Their number up to height $T$ is
\[
N(T)=\frac T{2\pi}\log\frac T{2\pi e}+\frac78+O(\log T)
\]
(Theorem~\ref{thm:rvm}). There are $O(\log T)$ of them in each window $]T,T+1]$ (Corollary~\ref{cor:nunit}), so near any point, $\zeta'/\zeta$ is the sum of $1/(s-\rho)$ over the nearby zeros, up to $O(\log|t|)$ (Proposition~\ref{prop:nearzeros}). And none lies to the right of $1-c/\log(|t|+2)$ (Theorem~\ref{thm:zerofree}). For large $|t|$ the proof gives the explicit gap $1/(14A'\log|t|)$. The third claim, that every non-trivial zero lies on $\operatorname{Re}s=\frac12$, is the Riemann hypothesis. The arguments of this chapter do not decide it.

Everything in Chapter~\ref{ch:explicit} rests on the chain from Jensen's formula to this zero-free region. That chapter makes Riemann's description of $\pi(x)$ precise: $\psi(x)$ equals $x$ minus a sum of waves $x^\rho/\rho$, one wave for each zero. With the zero-free region the waves are small, and the prime number theorem gets an error term. The explicit formula is the next question.

\section{Exercises}

\begin{exercise}[$\star$]\label{ex:digammahalf}
Differentiate the logarithm of Legendre's duplication formula (Volume~5, Theorem~5.14) to show that $\psi_\Gamma(\frac12)=-\gamma-2\log2$, and deduce that $B=\frac12\log4\pi-1-\frac\gamma2=-0.0230957\ldots$
\end{exercise}

\begin{exercise}[$\star\star$]\label{ex:bsum}
Show that $\sum_\rho\operatorname{Re}(1/\rho)$ converges absolutely and that $B=-\sum_\rho\operatorname{Re}(1/\rho)$. Deduce that if the Riemann hypothesis holds, then $\sum_\rho 1/|\rho|^2=-2B=0.0461914\ldots$
\end{exercise}

\begin{exercise}[$\star$]\label{ex:zerosums}
Show that $\sum_\rho|\rho|^{-1}$ diverges, while $\sum_\rho|\rho|^{-1}(\log|\rho|)^{-3}$ converges.
\end{exercise}

\begin{exercise}[$\star\star$, the signature exercise]\label{ex:nthzero}
Let $0<\gamma_1\le\gamma_2\le\cdots$ be the ordinates of the zeros in the upper half-plane. Show that $\gamma_n\sim\frac{2\pi n}{\log n}$, and that the average gap between consecutive ordinates near height $T$ is $\frac{2\pi}{\log(T/2\pi)}$. (At $T=10^6$ the average gap is $0.525$. Near the first zero the same formula gives $7.8$.)
\end{exercise}

\begin{exercise}[$\star\star$]\label{ex:carlson}
Let $f$ be entire, with $|f(z)|\le Ae^{c|z|}$ for some $c<\frac{\log2}2$, and $f(n)=0$ for every integer $n\ge1$. Show that $f=0$. What does $\sin\pi z$ say about the constant?
\end{exercise}

\begin{exercise}[$\star\star$]\label{ex:zetahalf}
Show that
\[
\zeta(s)=\sum_{n\le N}n^{-s}+\frac{N^{1-s}}{s-1}-s\int_N^\infty\{u\}u^{-s-1}\dd u
\]
for $\operatorname{Re}s>0$, $s\ne1$ and every integer $N\ge1$. Deduce that $\zeta(\frac12+it)=O(|t|^{1/2})$.
\end{exercise}

\begin{exercise}[$\star\star\star$]\label{ex:gammaratio}
Use \eqref{eq:stirlingvert} and the functional equation to show that $|\zeta(-1+it)|\asymp|t|^{3/2}|\zeta(2-it)|$ as $|t|\to\infty$, and more generally that
\[
|\zeta(\sigma+it)|=\Bigl(\frac{|t|}{2\pi}\Bigr)^{\frac12-\sigma}|\zeta(1-\sigma+it)|\Bigl(1+O\Bigl(\frac1{|t|}\Bigr)\Bigr)
\]
for fixed real $\sigma$.
\end{exercise}

\begin{exercise}[$\star$]\label{ex:xi2}
Using $\zeta(2)=\pi^2/6$ and $\Gamma(1)=1$, compute $\xi(2)$ from the definition in Volume~5, Theorem~10.8. Check that the value agrees with $\xi(-1)$, and that $-1$ is not a zero of $\xi$.
\end{exercise}

\begin{exercise}[$\star\star$]\label{ex:jensencheck}
\begin{enumerate}[label=(\alph*)]
\item Let $f(z)=e^z$ and let $R>0$. Check both sides of Jensen's formula on $|z|=R$.
\item Let $R>0$ and let $a\in\C$. If $a=0$, check the integral directly. If $a\ne0$ and $|a|\ne R$, show that
\[
\frac1{2\pi}\int_0^{2\pi}\log|Re^{i\theta}-a|\dd\theta=\log\max(R,|a|).
\]
\end{enumerate}
\end{exercise}

\begin{exercise}[$\star\star$]\label{ex:ordsum}
Let $N(T)$ count the zeros with positive ordinate at most $T$, as in Theorem~\ref{thm:rvm}. Show that
\[
\sum_{0<\gamma\le T}\frac1\gamma\sim\frac1{4\pi}(\log T)^2.
\]
\end{exercise}

\begin{exercise}[$\star\star\star$]\label{ex:fourteen}
\begin{enumerate}[label=(\alph*)]
\item In the proof of Theorem~\ref{thm:zerofree}, start from $4/(\sigma-\beta_0)\le3/(\sigma-1)+A'\log|t|$ and from the choice $\sigma=1+\delta/\log|t|$ with $\delta=1/(2A')$. Show where the factor $14$ in $1-\beta_0\ge1/(14A'\log|t|)$ comes from.
\item Fix $\sigma_0<1$ and a point $\sigma\in(1,2]$. Explain why the same inequality, at that fixed $\sigma$, does not forbid a zero with real part $\beta_0\le\sigma_0$.
\end{enumerate}
\end{exercise}

\section{Hints}

\begin{hint}{ex:digammahalf}
The duplication formula gives $\log\Gamma(x)+\log\Gamma(x+\frac12)=(1-2x)\log2+\frac12\log\pi+\log\Gamma(2x)$. Differentiate, then set $x=\frac12$, and use $\psi_\Gamma(1)=\Gamma'(1)=-\gamma$. Then compute $\xi'/\xi(1)=-B$.
\end{hint}
\begin{hint}{ex:bsum}
Use Theorem~\ref{thm:hadamard} at $s=1$, where $\xi'/\xi(1)=-B$, and the symmetry $\rho\mapsto1-\overline\rho$. Under the Riemann hypothesis, $\operatorname{Re}(1/\rho)=1/(2|\rho|^2)$.
\end{hint}
\begin{hint}{ex:zerosums}
Count the zeros with $2^k<|\rho|\le2^{k+1}$ by Theorem~\ref{thm:rvm}.
\end{hint}
\begin{hint}{ex:nthzero}
One has $N(\gamma_n)\approx n$. Invert $\frac T{2\pi}\log\frac T{2\pi}$ as in Exercise~\ref{ex:nthprime}.
\end{hint}
\begin{hint}{ex:carlson}
Apply Corollary~\ref{cor:jensen} on a disc centred at a point where $f\ne0$. A disc of radius $r$ contains about $r$ positive integers.
\end{hint}
\begin{hint}{ex:zetahalf}
Summation by parts on $\sum_{N<n\le X}n^{-s}$ for $\operatorname{Re}s>1$, then analytic continuation. Take $N=\lfloor|t|\rfloor$.
\end{hint}
\begin{hint}{ex:gammaratio}
Write the functional equation as
\[\zeta(s)=\pi^{s-\frac12}\frac{\Gamma((1-s)/2)}{\Gamma(s/2)}\zeta(1-s).\]
Compare $|\Gamma(\frac{1-\sigma}2-\frac{it}2)|$ and $|\Gamma(\frac\sigma2+\frac{it}2)|$ by \eqref{eq:stirlingvert}.
\end{hint}

\begin{hint}{ex:xi2}
The polynomial factor at $s=2$ is $\frac12\cdot2\cdot1=1$, and $\pi^{-1}\Gamma(1)\zeta(2)=\pi/6$. The functional equation moves the value to $-1$.
\end{hint}
\begin{hint}{ex:jensencheck}
(a) The real part of $Re^{i\theta}$ averages to $0$. (b) If $a=0$ the integrand is constant. If $0<|a|<R$, Jensen applies to $z-a$. If $|a|>R$, the disc contains no zero, and $\log|f(0)|=\log|a|$.
\end{hint}
\begin{hint}{ex:ordsum}
Write the sum as a Stieltjes integral against $N$, and integrate by parts. The main term of $N(t)$ is $(t/2\pi)\log(t/(2\pi e))$.
\end{hint}
\begin{hint}{ex:fourteen}
(a) Substitute $\delta=1/(2A')$ into $3/\delta+A'$, and subtract $\sigma-1=\delta/\log|t|$ from the resulting gap $\sigma-\beta_0$. (b) The right side grows like $A'\log|t|$, while a zero stuck at $\sigma_0$ would make the left side stay bounded.
\end{hint}

\section{Solutions}

\begin{solution}{ex:digammahalf}
\idea A known value of $\psi_\Gamma$ comes from an identity for $\Gamma$. Then $B$ is $-\xi'/\xi(1)$, computed factor by factor.
\writeup Differentiate the logarithm of the duplication formula:
\[
\psi_\Gamma(x)+\psi_\Gamma\bigl(x+\tfrac12\bigr)=-2\log2+2\psi_\Gamma(2x).
\]
At $x=\frac12$,
\[
\psi_\Gamma\bigl(\tfrac12\bigr)+\psi_\Gamma(1)=-2\log2+2\psi_\Gamma(1),
\]
so $\psi_\Gamma(\frac12)=\psi_\Gamma(1)-2\log2=-\gamma-2\log2$. Volume~5, Exercise~5.7, gives $\Gamma'(1)=-\gamma$.

From $\xi(s)=\xi(1-s)$ one has $\xi'(0)=-\xi'(1)$, so $B=-\xi'/\xi(1)$. Near $s=1$,
\[
\frac{\xi'}{\xi}(s)=\frac1s+\Bigl(\frac1{s-1}+\frac{\zeta'}{\zeta}(s)\Bigr)-\frac12\log\pi+\frac12\psi_\Gamma\Bigl(\frac s2\Bigr),
\]
and $1/(s-1)+\zeta'/\zeta(s)\to\gamma$. With $Z(s)=(s-1)\zeta(s)$ as in the proof of Theorem~\ref{thm:pnt}, one has $\zeta'/\zeta=Z'/Z-1/(s-1)$, and Volume~5, Proposition~10.6, gives $Z(1)=1$ and $Z'(1)=\gamma$. So
\[
\begin{aligned}\frac{\xi'}{\xi}(1)&=1+\gamma-\frac12\log\pi-\frac\gamma2-\log2\\&=1+\frac\gamma2-\frac12\log4\pi,\end{aligned}
\]
and $B=\frac12\log4\pi-1-\frac\gamma2$. Numerically, $\frac12\log(4\pi)=1.265512$, $\gamma/2=0.288608$, and $B=-0.0230957$.
\end{solution}

\begin{solution}{ex:bsum}
\idea The partial fractions at $s=1$ involve $1/(1-\rho)+1/\rho$. The symmetry of the zeros turns $1/(1-\rho)$ into another reciprocal, with the same real part.
\writeup One has $\operatorname{Re}(1/\rho)=\beta/|\rho|^2\in\bigl(0,1/|\rho|^2\bigr]$, so the series of real parts converges absolutely. Theorem~\ref{thm:hadamard} at $s=1$ gives
\[
-B=\frac{\xi'}{\xi}(1)=B+\sum_\rho\Bigl(\frac1{1-\rho}+\frac1\rho\Bigr),
\]
the series converging absolutely. Take real parts. The map $\rho\mapsto\rho^*=1-\overline\rho$ permutes the zeros (Proposition~\ref{prop:xi}(b)), and $1-\rho=\overline{\rho^*}$, so $\operatorname{Re}1/(1-\rho)=\operatorname{Re}(1/\rho^*)$. Summing a convergent series of positive terms in another order does not change the sum, and $\sum_\rho\operatorname{Re}1/(1-\rho)=\sum_\rho\operatorname{Re}(1/\rho)$. Since $B$ is real, $-2B=2\sum_\rho\operatorname{Re}(1/\rho)$. Under the Riemann hypothesis, $\beta=\frac12$, so
\[
\sum|\rho|^{-2}=2\sum\operatorname{Re}\frac1\rho=-2B=0.0461914.
\]
\end{solution}

\begin{solution}{ex:zerosums}
\idea There are about $(T/2\pi)\log T$ zeros up to height $T$. Each dyadic block $2^k<|\rho|\le2^{k+1}$ holds about $k\,2^k$ of them, each of size about $2^k$.
\writeup The modulus $|\rho|$ and the ordinate $|\gamma|$ differ by less than $1$, and zeros come in conjugate pairs. The number of zeros with $2^k<|\rho|\le2^{k+1}$ is
\[
2\bigl(N(2^{k+1})-N(2^k)\bigr)+O(\log2^k).
\]
Theorem~\ref{thm:rvm} says this is $(2^k/\pi)(k\log2+O(1))$, hence between $c_1k\,2^k$ and $c_2k\,2^k$ for large $k$. In this block $|\rho|^{-1}\ge2^{-k-1}$, so the block contributes at least $c_1k/2$ to $\sum|\rho|^{-1}$. The series diverges. With the extra factor $(\log|\rho|)^{-3}\le(k\log2)^{-3}$, the block contributes $O(k^{-2})$, and that series converges.
\end{solution}

\begin{solution}{ex:nthzero}
\idea One has $N(T)\approx(T/2\pi)\log T$. Inverting gives $T\approx2\pi n/\log n$, the same inversion as $p_n\approx n\log n$, with multiplication and division by $\log$ exchanged.
\writeup For every small $\varepsilon>0$, $N(\gamma_n)\ge n$ and $N(\gamma_n-\varepsilon)<n$. The jump of $N$ on an interval of length $1$ is $O(\log\gamma_n)$ by Corollary~\ref{cor:nunit}, so $N(\gamma_n)=n+O(\log\gamma_n)$. Theorem~\ref{thm:rvm} gives $(\gamma_n/2\pi)\log\gamma_n\sim n$. Taking logarithms, $\log\gamma_n\sim\log n$, and therefore $\gamma_n\sim2\pi n/\log n$.

The number of ordinates in $]T,T+H]$, for $H\to\infty$ with $H=o(T)$, is
\[
\frac H{2\pi}\log\frac T{2\pi}+O(\log T),
\]
because the derivative of the main term is $\frac1{2\pi}\log(T/(2\pi))$. The average gap is $2\pi/\log(T/2\pi)$. At $T=10^6$, $\log(10^6/(2\pi))=11.98$, and the gap is $0.525$. Near the first zero, $\log(14/(2\pi))=0.80$, and the formula gives $7.8$. Edwards records the first four positive ordinates as $14.13$, $21.02$, $25.01$ and $30.42$ (Volume~5 states the first three in the caption of its picture of the $s$-plane). The gaps between them are $6.9$, $4.0$ and $5.4$. Those numbers are taken from the tables. They are not recomputed here.
\end{solution}

\begin{solution}{ex:carlson}
\idea Zeros are expensive for a slowly growing function. The positive integers are evenly spread, about $r$ of them in a disc of radius $r$, while the budget for a disc of radius $2r$ is about $2cr$.
\writeup Suppose $f\ne0$. Replacing $f(z)$ by $f(z)/z^m$ for a suitable $m$ removes a zero at the origin and does not change the growth bound, up to the constant $A$, on $|z|\ge1$. The positive integers remain zeros. So assume $f(0)\ne0$. The disc $|z|\le r$ contains the $\lfloor r\rfloor$ zeros $1,\dots,\lfloor r\rfloor$. Corollary~\ref{cor:jensen} with $R=2r$ gives
\[
\lfloor r\rfloor\log2\le\log\frac{Ae^{2cr}}{|f(0)|}=2cr+O(1).
\]
Divide by $r$ and let $r\to\infty$. Then $\log2\le2c$, a contradiction.

Changing the ratio $R/r$ improves the constant. Carlson's theorem (1914), quoted from Titchmarsh, \emph{The Theory of Functions}, gives the sharp result, with growth $e^{c|z|}$ for any $c<\pi$ along the imaginary axis. The function $\sin\pi z$ vanishes at every integer and satisfies $|\sin\pi z|\le e^{\pi|z|}$, so some bound of this shape is necessary.
\end{solution}

\begin{solution}{ex:zetahalf}
\idea Summation by parts writes the tail of $\sum n^{-s}$ as an integral minus a fluctuating term. The formula then extends to $\operatorname{Re}s>0$. Cutting at $N\approx|t|$ balances the two pieces.
\writeup For $\operatorname{Re}s>1$ and $X>N$, Lemma~\ref{lem:abelint} with $a_n=1$ for $n>N$, $A(u)=\lfloor u\rfloor-N$ and $f(u)=u^{-s}$ gives
\[
\sum_{N<n\le X}n^{-s}=(\lfloor X\rfloor-N)X^{-s}+s\int_N^X(\lfloor u\rfloor-N)u^{-s-1}\dd u.
\]
Write $\lfloor u\rfloor=u-\{u\}$ and let $X\to\infty$:
\begin{align*}
\sum_{n>N}n^{-s}
&=s\int_N^\infty u^{-s}\dd u-Ns\int_N^\infty u^{-s-1}\dd u-s\int_N^\infty\{u\}u^{-s-1}\dd u\\
&=\frac{N^{1-s}}{s-1}-s\int_N^\infty\{u\}u^{-s-1}\dd u.
\end{align*}
The right side is holomorphic on $\operatorname{Re}s>0$ except at $s=1$. The integral is bounded by $\int_N^\infty u^{-\sigma-1}\mathrm du$, uniformly on compact sets. Split the integral at the integers and apply Volume~5, Theorem~9.15. The identity extends.

At $s=\frac12+it$ with $|t|\ge2$ and $N=\lfloor|t|\rfloor$, one has $|\sum_{n\le N}n^{-1/2}|\le2\sqrt N$, $|N^{1/2-it}/(s-1)|\le\sqrt N/|t|\le1$, and
\[
\Bigl|s\int_N^\infty\{u\}u^{-3/2-it}\dd u\Bigr|\le|s|\cdot\frac2{\sqrt N}=O(\sqrt{|t|}).
\]
So $\zeta(\frac12+it)=O(|t|^{1/2})$. This is far from the truth. The improvements are quoted from Titchmarsh, Chapter~V. Convexity gives $|t|^{1/4+\varepsilon}$. Weyl's exponential sums, in the hands of Hardy and Littlewood, give $|t|^{1/6+\varepsilon}$. Bourgain (2017) reached $|t|^{13/84+\varepsilon}$. The Lindelöf hypothesis, which follows from the Riemann hypothesis, predicts $|t|^\varepsilon$.
\end{solution}

\begin{solution}{ex:gammaratio}
\idea The functional equation relates $\zeta$ at $s$ and at $1-s$ through a ratio of Gamma values. On vertical lines, Stirling says that the exponentials $e^{-\pi|t|/4}$ cancel, and a power of $|t|$ survives.
\writeup From $\xi(s)=\xi(1-s)$,
\[
\pi^{-s/2}\Gamma\Bigl(\frac s2\Bigr)\zeta(s)=\pi^{-(1-s)/2}\Gamma\Bigl(\frac{1-s}2\Bigr)\zeta(1-s),
\]
which is the formula in the hint. Let $s=\sigma+it$ with $t\to+\infty$. By \eqref{eq:stirlingvert},
\[
\log|\Gamma(a+ib)|=\bigl(a-\tfrac12\bigr)\log b-\frac{\pi b}2+\tfrac12\log(2\pi)+O\Bigl(\frac1b\Bigr)
\]
for fixed $a$. Also $|\Gamma(\overline z)|=|\Gamma(z)|$. With $b=t/2$,
\[
\begin{aligned}\log\Bigl|\frac{\Gamma(\frac{1-\sigma}2-\frac{it}2)}{\Gamma(\frac\sigma2+\frac{it}2)}\Bigr|
&=\Bigl(\frac{1-\sigma}2-\frac\sigma2\Bigr)\log\frac t2+O\Bigl(\frac1t\Bigr)\\
&=\Bigl(\frac12-\sigma\Bigr)\log\frac t2+O\Bigl(\frac1t\Bigr).\end{aligned}
\]
And $|\pi^{s-1/2}|=\pi^{\sigma-1/2}$. So
\[
|\zeta(\sigma+it)|=\Bigl(\frac t{2\pi}\Bigr)^{\frac12-\sigma}|\zeta(1-\sigma-it)|\Bigl(1+O\Bigl(\frac1t\Bigr)\Bigr),
\]
and $|\zeta(1-\sigma-it)|=|\zeta(1-\sigma+it)|$ by conjugation. For $\sigma=-1$,
\[
|\zeta(-1+it)|\sim\Bigl(\frac t{2\pi}\Bigr)^{3/2}|\zeta(2+it)|,
\]
and $|\zeta(2+it)|$ lies between $\zeta(4)/\zeta(2)$ and $\zeta(2)$. To the left of the strip, $\zeta$ grows like a power of $|t|$.
\end{solution}

\begin{solution}{ex:xi2}
\idea At $s=2$ every factor in the definition is elementary. The functional equation carries the value to $s=-1$, where $\zeta$ itself is not being evaluated by the series.
\writeup Volume~5, Theorem~10.8, gives
\[
\xi(2)=\frac12\cdot2\cdot1\cdot\pi^{-1}\cdot\Gamma(1)\cdot\zeta(2)=\pi^{-1}\cdot\frac{\pi^2}6=\frac\pi6.
\]
The same theorem gives $\xi(-1)=\xi(2)=\pi/6\ne0$. So $-1$ is not a zero. The trivial zeros of $\zeta$ are at $-2,-4,-6,\dots$, where the pole of $\Gamma(s/2)$ cancels a zero of the polynomial factor in a way that leaves $\xi$ non-zero. At $s=-1$, $\Gamma(s/2)=\Gamma(-1/2)$ is finite.
\end{solution}

\begin{solution}{ex:jensencheck}
\idea The exponential has no zeros, so Jensen says that the average of its real part on a circle is the value at the centre. A single zero at $a$ produces the mean of $\log|z-a|$.
\writeup (a) The function $e^z$ has no zeros, and $|e^0|=1$, so the left side of Jensen's formula is $0$. On $|z|=R$,
\[
\log|e^{Re^{i\theta}}|=R\cos\theta,
\]
and the average of $\cos\theta$ over $[0,2\pi]$ is $0$.

(b) If $a=0$, the integrand is the constant $\log R$, and $\max(R,0)$ is not the right comparison: the integral equals $\log R$ by inspection. Now let $a\ne0$ and $|a|\ne R$. If $|a|<R$, the function $f(z)=z-a$ has one zero inside the circle, at $a$, and $f(0)=-a\ne0$. Jensen's formula says
\[
\log|a|+\log\frac R{|a|}=\frac1{2\pi}\int_0^{2\pi}\log|Re^{i\theta}-a|\dd\theta.
\]
The left side is $\log R$. If $|a|>R$, the same $f$ has no zero in $|z|\le R$, so the formula says that the average equals $\log|f(0)|=\log|a|$. In both cases the average equals $\log\max(R,|a|)$.
\end{solution}

\begin{solution}{ex:ordsum}
\idea Partial summation turns a sum of $1/\gamma$ into an integral of $N(t)/t^2$. The main term of $N$ is a multiple of $t\log t$, and $\int(\log t)/t\dd t$ is $\frac12(\log t)^2$.
\writeup Let $\gamma$ run through the positive ordinates, repeated by multiplicity, and write $N$ for the counting function of Theorem~\ref{thm:rvm}. For $T$ larger than the first ordinate,
\[
\sum_{0<\gamma\le T}\frac1\gamma=\int_{\gamma_1^-}^T\frac1t\,\mathrm dN(t)=\frac{N(T)}T+\int_{\gamma_1}^T\frac{N(t)}{t^2}\dd t.
\]
The boundary term is $O(\log T)$, because $N(T)=O(T\log T)$. Insert
\[
N(t)=\frac t{2\pi}\log\frac t{2\pi e}+\frac78+O(\log t).
\]
Then
\[
\int_2^T\frac1{t^2}\cdot\frac t{2\pi}\log\frac t{2\pi e}\dd t=\frac1{2\pi}\int_2^T\frac1t\log\frac t{2\pi e}\dd t=\frac1{4\pi}\Bigl(\log\frac T{2\pi e}\Bigr)^2+O(1).
\]
The integral of $(7/8)/t^2$ converges as $T\to\infty$, and $\int_2^T(\log t)/t^2\dd t$ stays bounded. So
\[
\sum_{0<\gamma\le T}\frac1\gamma=\frac1{4\pi}(\log T)^2+O(\log T),
\]
which is asymptotic to $(\log T)^2/(4\pi)$.
\end{solution}

\begin{solution}{ex:fourteen}
\idea The pole contributes a $3$, and a zero on the same vertical line would contribute a $4$. The surplus $4-3$ is what pushes $\beta$ to the left of $1$. Letting $\sigma$ slide towards $1$ at speed $1/\log|t|$ balances the pole against the $O(\log|t|)$ error, and the arithmetic produces $14$.
\writeup (a) With $\sigma=1+\delta/\log|t|$ and $\delta=1/(2A')$,
\[
\frac3{\sigma-1}+A'\log|t|=\Bigl(\frac3\delta+A'\Bigr)\log|t|.
\]
Since $1/\delta=2A'$, one has $3/\delta=6A'$, and $3/\delta+A'=7A'$. The inequality becomes $4/(\sigma-\beta_0)\le7A'\log|t|$, so
\[
\sigma-\beta_0\ge\frac4{7A'\log|t|}.
\]
Subtract $\sigma-1=\delta/\log|t|=1/(2A'\log|t|)$:
\[
1-\beta_0\ge\Bigl(\frac4{7A'}-\frac1{2A'}\Bigr)\frac1{\log|t|}.
\]
The coefficient is $(8-7)/(14A')=1/(14A')$, because $4/7-1/2=(8-7)/14$.

(b) Fix $\sigma\in(1,2]$ and suppose $\beta_0\le\sigma_0<1$. Then $\sigma-\beta_0\ge\sigma-\sigma_0$, so
\[
\frac4{\sigma-\beta_0}\le\frac4{\sigma-\sigma_0}.
\]
The $3$-$4$-$1$ inequality only says that the left side is at most $3/(\sigma-1)+A'\log|t|$. For large $|t|$ that upper bound tends to infinity, and a bounded left side sits under it. A zero parked at real part $\sigma_0$ produces no contradiction at this fixed $\sigma$. The proof reaches a shrinking region because $\sigma$ itself moves toward $1$. Then $3/(\sigma-1)$ grows like $\log|t|$, and the surplus of $4$ over $3$ leaves a gap of size $1/(14A'\log|t|)$.
\end{solution}

\section*{Sources and further reading}

\begin{itemize}
\item B.~Riemann, \emph{Ueber die Anzahl der Primzahlen unter einer gegebenen Grösse}, Monatsberichte der Berliner Akademie (1859), 671--680. The English translation in H.~M.~Edwards's book, below, is the source of the epigraph and of the phrase ``fleeting vain attempts''. Riemann's $\xi(t)$ is our $\xi(\frac12+it)$.
\item H.~M.~Edwards, \emph{Riemann's Zeta Function}, Academic Press, 1974 (Dover reprint). A line-by-line commentary on the memoir, and tables of the early ordinates.
\item J.~Hadamard, \emph{Étude sur les propriétés des fonctions entières et en particulier d'une fonction considérée par Riemann}, Journal de Mathématiques Pures et Appliquées 9 (1893), 171--215.
\item H.~von Mangoldt, \emph{Zur Verteilung der Nullstellen der Riemannschen Funktion $\xi(t)$}, Mathematische Annalen 60 (1905), 1--19; C.-J.~de la Vallée Poussin, \emph{Sur la fonction $\zeta(s)$ de Riemann et le nombre des nombres premiers inférieurs à une limite donnée}, Mémoires couronnés de l'Académie royale de Belgique 59 (1899).
\item J.~L.~W.~V.~Jensen, \emph{Sur un nouvel et important théorème de la théorie des fonctions}, Acta Mathematica 22 (1899), 359--364. For Jensen's work at the Copenhagen telephone company, see the MacTutor biography.
\item H.~Davenport, \emph{Multiplicative Number Theory}, 3rd ed., Springer GTM~74, 2000, Chapters~11--15.
\item H.~L.~Montgomery and R.~C.~Vaughan, \emph{Multiplicative Number Theory I}, Cambridge, 2007, Chapters~6 and~10. The local lemma is their Lemma~6.3.
\item E.~C.~Titchmarsh, \emph{The Theory of the Riemann Zeta-Function}, 2nd ed., revised by D.~R.~Heath-Brown, Oxford, 1986, Chapters~9 and~10, and Chapter~VI for the Vinogradov--Korobov region.
\item A.~M.~Turing, \emph{Some calculations of the Riemann zeta-function}, Proceedings of the London Mathematical Society 3 (1953), 99--117; X.~Gourdon, \emph{The $10^{13}$ first zeros of the Riemann zeta function, and zeros computation at very large height}, preprint, 2004.
\item Zeros on the line, quoted and not proved here: G.~H.~Hardy, \emph{Sur les zéros de la fonction $\zeta(s)$ de Riemann}, Comptes Rendus 158 (1914), 1012--1014; A.~Selberg, \emph{On the zeros of Riemann's zeta-function}, Skrifter Norske Videnskaps-Akademi Oslo I, 1942, no.~10; N.~Levinson, \emph{More than one third of zeros of Riemann's zeta-function are on $\sigma=1/2$}, Advances in Mathematics 13 (1974), 383--436; J.~B.~Conrey, \emph{More than two fifths of the zeros of the Riemann zeta function are on the critical line}, Journal für die reine und angewandte Mathematik 399 (1989), 1--26; K.~Pratt, N.~Robles, A.~Zaharescu and D.~Zeindler, \emph{More than five-twelfths of the zeros of $\zeta$ are on the critical line}, Research in the Mathematical Sciences 7 (2020), article~2.
\item J.~Bourgain, \emph{Decoupling, exponential sums and the Riemann zeta function}, Journal of the American Mathematical Society 30 (2017), 205--224; E.~C.~Titchmarsh, \emph{The Theory of Functions}, 2nd ed., Oxford, 1939, Chapter~V, for Carlson's theorem.
\end{itemize}
